% -----------------------------------------------------------------------------
% Conteudo AODV-EN. Template base: Prof. Wesley Pacheco Calixto (IFG).
% -----------------------------------------------------------------------------
\addcontentsline{toc}{section}{Abstract}

\section*{Abstract}

Low-cost wireless mesh networks built on top of the ESP-NOW protocol lack native support
for multi-hop routing, which limits their range and scalability. This work proposes,
implements, and evaluates AODV-EN, an adaptation of the Ad-hoc On-Demand Distance Vector
protocol (AODV, RFC 3561) to operate over ESP-NOW on ESP32 modules, featuring MAC-based
addressing, route request dissemination through controlled flooding over link-layer
broadcast, route maintenance with sequence numbers, and end-to-end acknowledgment.
Following the Design Science Research paradigm, the protocol core was implemented in C as
a reusable component and compared against a flooding-based reference algorithm, measuring
Packet Delivery Ratio (PDR), end-to-end latency, Normalized Routing Load (NRL), and
estimated energy consumption. In the hardware scenario C1 (ten ESP32 nodes in a
multi-hop chain of three to four hops, thirty repetitions per algorithm), AODV-EN
delivered 94.1\% of packets versus only 13.9\% for flooding --- a striking gap in the
multi-hop regime, where flooding collapses for lack of a route and retransmission ---
while costing about 20\% fewer transmissions per delivered packet (54.2 versus 68.8). A
second hardware scenario (C3, a dense ten-node mesh, five hops) showed that this advantage
is topology-dependent: in the mesh, path redundancy raises flooding's delivery to 77.8\%,
narrowing --- without eliminating --- AODV-EN's advantage (91.4\%). Simulation analysis
confirmed that the channel-occupancy disadvantage of flooding grows with network scale,
with AODV-EN's efficiency advantage consolidating beyond about eleven nodes. A third hardware scenario (C4, six
nodes), with cyclic failure of an on-path relay, demonstrated the protocol's self-healing:
AODV-EN sustained 70.6\% delivery by rediscovering the route through an alternative relay
after each outage. Regarding estimated energy efficiency,
although AODV-EN's raw aggregate consumption is higher, its energy per successfully
delivered packet is about five times lower than flooding's (0.69 versus 3.75~J/packet),
since routing delivers far more. It is concluded that AODV-EN's reactive routing becomes
advantageous as the network grows in size and diameter, enabling standardized, low-cost
mesh networks over ESP-NOW that are energy-efficient per useful delivery.

\vspace{1cm}

\noindent
\textbf{Keywords:} mesh networks; ESP-NOW; ESP32; AODV; ad-hoc routing;
Internet of Things.
